\documentclass[12pt]{article}
\usepackage[utf8]{inputenc}
\usepackage[T1]{fontenc}
\usepackage{hyperref}
\title{Tongan Language Resources for K-12 Curriculum Development}
\author{Compiled by Claude Sonnet 5 (claude-sonnet-5) for Andr\'as Kornai}
\date{}
\begin{document}
\maketitle

\section{Introduction}
Tongan (lea faka-Tonga) is the national language of Tonga and is taught in schools there. It is also taught in New Zealand. This report is based on one search round. No Tongan Ministry of Education syllabus or Tongan-language dictionary appeared in the results.

\section{Primary Educational Resources}
\begin{itemize}
\item New Zealand Ministry of Education, \emph{Ko e Fakahinohino ki he Lea Faka-Tonga: The Tongan Language Guidelines}, aligned with the New Zealand Curriculum for Years 1--13: \url{https://newzealandcurriculum.tahurangi.education.govt.nz/the-tongan-language-guidelines/5637168581.p}
\item Teaching and learning lea faka Tonga (NZ): \url{https://newzealandcurriculum.tahurangi.education.govt.nz/teaching-and-learning-lea-faka-tonga/5637175354.p}
\end{itemize}
The NZ materials mention NCEA levels 1--3 Tongan achievement standards.

\section{Linguistic and Reference Materials}
The NZ guidelines include sections on resources such as grammars, dictionaries and coursebooks, but I did not see the individual titles.

\section{Religious and Community Resources}
None found in this search.

\section{Digital and Audio Resources}
None found in this search.

\section{Curriculum Development Context}
A 2012 Matangi Tonga report says Tongan is compulsory from Class 1 through Form 7 under the Tongan language policy, with English compulsory from Class 3 to Form 6: \url{https://matangitonga.to/node/7201}. A 2007 curriculum reform contracted to the University of Queensland (UniQuest) is mentioned, and I did not find its outcome. The learner.net pages found were unreliable and are not cited.
\end{document}